\bookchapter{Chapitre 31 : Initiation au tournage de vidéos courtes et à la création de contenu}

\chapterintro{En 2026, les courtes vidéos ne sont plus « un jouet pour les jeunes », elles sont devenues le moyen de base pour tous d'obtenir des informations et de s'exprimer. Que vous souhaitiez enregistrer votre vie, partager vos connaissances ou gagner un revenu complémentaire, le smartphone est votre meilleur outil de création — pas besoin d'acheter un appareil photo ni d'apprendre Premiere, un smartphone suffit pour commencer. Ce chapitre vous aide à démarrer de zéro.}

\booksection{I. Avant de filmer, réfléchissez : pour qui filmez-vous ?}

La première erreur de beaucoup de gens est de « prendre le smartphone et filmer » sans savoir pour qui. Avant de tourner, répondez à trois questions :

\begin{enumerate}
  \item \textbf{Qui est mon audience ?}— Filmer pour des enfants, pour des pairs ou pour des seniors implique des styles de contenu totalement différents.
  \item \textbf{Quelle valeur puis-je apporter ?}— Faire rire (divertissement), faire apprendre (connaissances), faire saliver (gastronomie), toucher (émotions). Au moins un de ces éléments.
  \item \textbf{Sur quelle plateforme puis-je publier régulièrement ?}— Douyin/Kuaishou (divertissement general), Xiaohongshu (Mode de vie), Bilibili (contenu approfondi), Chaînes vidéo (cercle social). Choisissez une plateforme principale et approfondissez, ne vous dispersez pas sur toutes les plateformes dès le départ.
\end{enumerate}

\booksection{II. Réglages de base du smartphone — 5 choses à faire avant de filmer}

\begin{enumerate}
  \item \textbf{Résolution 1080p, fréquence d'images 30 fps}: C'est le standard des plateformes de courtes vidéos. Les fichiers 4K sont trop volumineux et lents à télécharger, 60 fps n'offre aucun avantage sur les plateformes sociales et produit des fichiers plus volumineux.
  \item \textbf{Nettoyer l'objectif}: La première et la plus importante. Essuyez l'objectif avec votre vêtement avant de filmer, l'image sera immédiatement plus nette.
  \item \textbf{Verrouiller la mise au point et l'exposition}: Dans le viseur, appuyez longuement (2 secondes) sur le sujet, le message « Verrouillage AE/AF » apparaîtra. Ainsi, l'image ne deviendra pas soudainement claire ou sombre pendant le tournage.
  \item \textbf{Horizontalement ou verticalement ?}— Pour Douyin/Kuaishou/Chaînes vidéo/Xiaohongshu\textbf{filmez verticalement}(9:16) ; pour Bilibili/YouTube\textbf{filmez horizontalement}(16:9). Définissez la plateforme avant de choisir l'orientation.
  \item \textbf{Activer la grille}: Réglages > Appareil photo > Grille. La base de la règle des tiers, pour trouver rapidement le centre de l'image.
\end{enumerate}

\booksection{III. Cinq mouvements de caméra — faites « bouger » l'image}

Le plus grand problème des débutants est une image trop « statique » — filmer immobile du début à la fin. Maîtrisez ces cinq mouvements de base, l'image aura immédiatement un aspect professionnel :

\begin{itemize}
  \item \textbf{Pousser (Push)}: Le smartphone se rapproche lentement du sujet. Crée un sentiment de tension ou d'intimité.
  \item \textbf{Tirer (Pull)}: Le smartphone s'éloigne lentement du sujet. Du détail au plan d'ensemble, montre l'environnement.
  \item \textbf{Panorama (Pan)}: Le smartphone reste sur place, pivote à gauche/droite. Idéal pour les paysages et la présentation d'espaces.
  \item \textbf{Suivre (Track)}: Le smartphone se déplace parallèlement, suivant le sujet. Nécessite une tenue stable, idéalement avec un stabilisateur.
  \item \textbf{Accompagner (Follow)}: Le smartphone suit le sujet par derrière ou par devant. Le mouvement le plus courant en vlog.
\end{itemize}

\textbf{La stabilité avant tout}: Les images tremblotantes semblent peu professionnelles. Tenir à deux mains est plus stable qu'à une seule ; en marchant, pliez légèrement les genoux, utilisez la « marche ninja » pour réduire les secousses ; si vous en avez le budget, achetez un stabilisateur (DJI OM, environ 500-800 yuans), l'amélioration est qualitative.

\booksection{IV. Introduction au montage — Jianying (CapCut) suffit}

Jianying (CapCut, version chinoise de CapCut) est l'outil de montage sur smartphone le plus adapté aux débutants : gratuit, complet et riche en modèles. Voici le flux de travail de base :

\begin{enumerate}
  \item \textbf{Importer les rushs}: Ouvrez Jianying (CapCut) > Commencer la création > Sélectionnez les clips vidéo filmés et importez-les.
  \item \textbf{Couper}: Faites glisser la timeline, coupez les parties inutiles au début et à la fin de chaque clip. Principe : ne gardez que les parties informatives, mieux vaut court que long.
  \item \textbf{Ajouter des transitions}: Entre deux clips, touchez le bouton « Transition », choisissez un effet simple. Les débutants devraient utiliser « Fondu » ou « Fondre au noir », pas de transitions 3D voyantes.
  \item \textbf{Ajouter des sous-titres}: Jianying (CapCut) dispose d'une fonction « reconnaissance de sous-titres », générant automatiquement les sous-titres en un seul toucher. Précision d'environ 90 \%, une relecture manuelle suffit.
  \item \textbf{Ajouter de la musique}: La bibliothèque musicale de Jianying (CapCut) contient de nombreuses musiques libres de droits. Choisissez un BGM dont le rythme correspond au contenu, réglez le volume à environ 1/3 du son original de la vidéo.
  \item \textbf{Exporter}: Exportez en 1080p, 30 fps, débit « recommandé ».
\end{enumerate}

\booksection{V. L'importance de la couverture — 70 \% des gens ne regardent que la couverture pour décider}

La couverture d'une courte vidéo est la « vitrine », elle détermine si les gens vont cliquer. Quelques points clés :

\begin{itemize}
  \item \textbf{La couverture doit avoir une phrase}: Écrivez en gros le contenu principal de la vidéo, par exemple « Apprenez la retouche photo sur smartphone en 3 minutes ». Police grande, couleur vive, contraste avec le fond.
  \item \textbf{La couverture et le contenu doivent correspondre}: Ne faites pas de « putaclic » — si la couverture dit une chose et le contenu une autre, les spectateurs zapperont immédiatement et l'algorithme pénalisera la vidéo.
  \item \textbf{Les expressions doivent être vivantes}: S'il y a une personne sur la couverture, choisissez un instant d'expression vive ou exagérée. Les visages attirent plus de clics que le texte.
  \item \textbf{Faire une couverture avec Jianying}: Dans Jianying (CapCut), sélectionnez la meilleure image, ajoutez du texte superposé, exportez en image comme couverture.
\end{itemize}

\booksection{VI. Trois principes de création de contenu}

\begin{itemize}
  \item \textbf{Les 3 premières secondes sont cruciales}: Sur les plateformes de courtes vidéos, les spectateurs décident en 3 secondes de zapper ou de continuer. Les 3 premières secondes doivent offrir un « crochet » — une question, un suspense, une image inattendue.
  \item \textbf{Une vidéo ne traite qu'un seul sujet}: L'erreur la plus courante des débutants est de mettre trop de contenu dans une vidéo. Une vidéo ne résout qu'un seul problème ; mieux vaut traiter à fond que survoler.
  \item \textbf{La persévérance compte plus que le talent}: La recommandation algorithmique nécessite de « nourrir le compte ». Les 20 premières vidéos n'auront presque pas de vues, c'est normal. Persévérez avec 2-3 publications par semaine pour que l'algorithme comprenne votre direction éditoriale, le trafic augmentera progressivement.
\end{itemize}

\begin{tcbpractice}{Exercice : filmez une courte vidéo de 30 secondes}
\begin{enumerate}
  \item Choisissez un petit sujet que vous maîtrisez (par exemple « comment agrandir la police du smartphone »).
  \item Filmez 3-5 clips en mode vertical avec le smartphone, 5-10 secondes chacun.
  \item Utilisez Jianying (CapCut) pour assembler, couper, ajouter des sous-titres et un BGM.
  \item Créez une couverture avec du texte en gros.
  \item Après export, envoyez d'abord à un ami pour recueillir des retours avant de publier.
\end{enumerate}

\end{tcbpractice}

\begin{tcbsummary}{Résumé du chapitre}
\begin{itemize}
  \item Réfléchissez avant de filmer : qui est l'audience, quelle valeur apporter, sur quelle plateforme publier
  \item 5 choses avant de filmer : 1080p/30 fps, nettoyer l'objectif, verrouiller mise au point et exposition, déterminer l'orientation, activer la grille
  \item Cinq mouvements : pousser, tirer, panorama, suivre, accompagner — la stabilité avant tout
  \item Montage avec Jianying (CapCut) : couper transitions sous-titres BGM exporter, en six étapes
  \item Les 3 premières secondes sont cruciales, une vidéo un seul sujet, la persévérance compte plus que le talent
\end{itemize}

\end{tcbsummary}
